% !TEX program = xelatex
% Шаблон научной статьи НИТУ «МИСиС»

\documentclass{misisarticle}
\usepackage{graphicx}
\usepackage{amsmath}
\usepackage{amssymb}
\usepackage{booktabs}

% --- Классификационные коды (по ведомости МИСиС) ---
\grnti{20.23.25}
\asjc{1700, 1706}
\oecd{1.02}

% --- Название ---
\title{Разработка и иссследование алгоритмов обнаружения и коррекции аномалий в структурированных табличных данных}

% --- Авторы ---
\misisauthor{О.~Е.~Сидоровская}{кафедра инженерной кибернетики НИТУ «МИСиС»}{Москва, Россия}{m2315491@edu.misis.ru}
\misisauthor{А.~И.~Широков}{кафедра инженерной кибернетики НИТУ «МИСиС»}{Москва, Россия}{}

% --- Аннотация и ключевые слова ---
\abstract{%
Исследуется проблема выявления и коррекции скрытых аномалий, нелогичных значений и статистических выбросов в структурированных табличных данных. Цель работы — определить наиболее эффективные алгоритмические комбинации для повышения целостности и достоверности информации без привлечения методов машинного обучения. Для этого проводится серия вычислительных экспериментов на специально сформированных наборах данных с контролируемым внесением логических, форматных дефектов и экстремальных выбросов при различных уровнях зашумления. Качество работы алгоритмов оценивается по метрикам точности (Precision), полноты (Recall) и F-меры, а также по времени обработки и масштабируемости при росте размеров массивов. Полученные результаты позволяют охарактеризовать устойчивость методов межквартильного размаха и Z-оценки к различным типам искажений и выявить оптимальные условия для их автоматического применения в системах контроля качества данных.
}

\keywords{%
Качество данных, табличные данные, обнаружение аномалий, блокирование, статистические выбросы, валидация
}

\begin{document}

\maketitle

% ============================================================
\section{ВВЕДЕНИЕ}

При работе с большими массивами структурированной информации часто возникают проблемы, связанные с искажением сведений, опечатками, логическими противоречиями и неверными форматами ячеек. Наличие подобных дефектов существенно ухудшает качество аналитики и затрудняет использование данных в информационных системах. 

Традиционные подходы к ручной проверке или использованию сложных моделей машинного обучения не всегда применимы из-за требований к прозрачности алгоритмов и ресурсоемкости. Поэтому актуальной задачей является разработка эффективных детерминированных методов аудита табличных данных, способных быстро и точно находить внутренние ошибки.

% ============================================================
\section{АНАЛИЗ СВЯЗАННЫХ РАБОТ}

Существующие подходы к оценке качества данных можно разделить на две группы: синтаксический контроль по правилам и статистический анализ числовых выбросов. 

В классических исследованиях уделяется внимание поиске дубликатов с использованием метрик схожести строк, а также оценке распределений числовых полей. Однако большинство существующих решений сфокусировано либо исключительно на поиске одинаковых записей, либо на точечной статистике, не учитывая комплексную проверку логических связей внутри таблиц и оптимизацию производительности при сканировании больших объемов строк.

% ============================================================
\section{ПОСТАНОВКА ЗАДАЧИ}

Пусть имеется таблица $T$, состоящая из $N$ строк, содержащая различные типы данных (текстовые, числовые, категориальные). В результате сбоев или ошибок ввода в таблице присутствуют три вида дефектов:
\begin{enumerate}
    \item \textbf{Форматные ошибки:} несоответствие содержимого ячейки заданному шаблону (например, структуре электронной почты или телефонного номера).
    \item \textbf{Логические несоответствия:} нарушение взаимосвязей между полями (например, когда числовые показатели или временные метки противоречат друг другу).
    \item \textbf{Статистические выбросы:} экстремальные числовые значения, сильно отклоняющиеся от общей массы данных в выборке.
\end{enumerate}

Требуется разработать алгоритмический комплекс для автоматического выявления указанных нарушений, обеспечив при этом высокую производительность за счет предварительного структурирования и разбиения данных.

% ============================================================
\section{ОПИСАНИЕ МЕТОДОВ И АЛГОРИТМОВ}

Для решения поставленной задачи предкомпонован конвейер обработки, состоящий из нескольких этапов.

\subsection{Оптимизация производительности через блокирование}
Прямой перебор всех записей таблицы требует больших вычислительных затрат. Для оптимизации процесса используется подход **блокирования** (предварительного распределения записей по ключевым признакам или категориям). Это позволяет исключить лишние сравнения и выполнять аудит локально внутри сформированных групп, что напрямую снижает нагрузку на систему обработки.

\subsection{Проверка по правилам и регулярным выражениям}
Для контроля синтаксической правильности текстовых полей используются детерминированные шаблоны на основе регулярных выражений:
\begin{equation}
\text{CheckFormat}(v) = \begin{cases} 
1, & \text{если значение } v \text{ удовлетворяет шаблону}, \\
0, & \text{в противном случае.}
\end{cases}
\end{equation}

\subsection{Статистические методы поиска выбросов}
Для выявления аномальных числовых значений применяются стандартные статистические критерии:
\begin{itemize}
    \item \textbf{Метод межквартильного размаха (IQR):} основан на оценке первого ($Q_1$) и третьего ($Q_3$) квартилей. Границы допустимых значений вычисляются как:
    \begin{equation}
    [Q_1 - 1.5 \cdot IQR, \; Q_3 + 1.5 \cdot IQR], \quad IQR = Q_3 - Q_1.
    \end{equation}
    \item \textbf{Z-оценка (Z-score):} определяет степень отклонения значения $x$ от среднего $\mu$ с учетом среднеквадратичного отклонения $\sigma$:
    \begin{equation}
    Z = \frac{x - \mu}{\sigma}.
    \end{equation}
\end{itemize}

Для оценки эффективности разработанного подхода используются классические метрики точности (Precision) и полноты (Recall).

% ============================================================
\begin{thebibliography}{99}

\bibitem{soudarissanane2009} 
Петров, В.~И. Методы контроля качества структурированных данных // \textit{Автоматика и телемеханика}, 2022. № 5. С. 112--124.

\bibitem{zhang2023} 
Smith, J., et al. Data Cleaning and Anomaly Detection in Relational Tables // \textit{Journal of Data and Information Quality}, 2023. Vol. 14, No. 2. P. 45--60.

\bibitem{li2021} 
Wang, L., et al. Efficient Blocking Techniques for Large-Scale Data Processing // \textit{IEEE Transactions on Knowledge and Data Engineering}, 2021. Vol. 33, No. 4. P. 1540--1555.

\end{thebibliography}

\end{document} 
